% LOMBEIDA ESCALERAS — solo contenido, sin preámbulo
\section*{LOMBEIDA ESCALERAS — Tolerancia a fallos (2 min) y lecciones/conclusiones (1 min)}
\paragraph{Diapositiva 10 (aprox. 1 min).}
Spark divide el trabajo en tareas que ejecutan distintos \textit{executors} \cite{sparkRddGuide}. Si una tarea falla, puede reintentarla: el límite de fallos consecutivos de esa tarea se configura con \texttt{spark.task.maxFailures} \cite{sparkConfiguration}. Cuando se pierde una partición de un RDD, Spark puede reconstruirla usando el linaje de transformaciones que la creó \cite{sparkRddGuide}. Si falta una salida de \textit{shuffle}, el fallo de lectura hace regresar a la etapa productora para regenerar los datos \cite{sparkFetchFailed}. Así, la recuperación depende de qué parte del trabajo paralelo se perdió \cite{sparkRddGuide,sparkFetchFailed}.

\paragraph{Diapositiva 11 (aprox. 1 min).}
El caso oficial SPARK-46182, registrado en 2023, describe una condición de carrera durante el retiro de un \textit{executor}; el informe no cuantifica usuarios afectados \cite{spark46182}. Primero, la comprobación de bloques \textit{shuffle} no encontró pendientes. Después terminó una tarea y produjo nuevos bloques. El executor recibió la orden de detenerse antes de que el hilo de migración detectara esos bloques. La salida se perdió y Spark tuvo que recalcular la partición \cite{spark46182}. El síntoma fue la pérdida de datos temporales; la causa inmediata, el retiro prematuro; la causa raíz, la falta de coordinación entre el fin de la tarea y la comprobación de migración \cite{spark46182}. El registro figura como corregido para las versiones 3.4.3, 3.5.1 y 4.0.0 \cite{spark46182}. Antes del ensayo, contrastar esta cadena con los campos exigidos por FORM-AC-01.

\paragraph{Diapositiva 12 (aprox. 1 min).}
Para el RTV propongo dos controles; no afirmo que ya estén implementados. Primero, confirmar cada resultado antes de permitir el siguiente avance. Las transacciones de PostgreSQL permiten agrupar escrituras para que se confirmen juntas o no tengan efecto si falla la operación \cite{postgresTransactions}. Segundo, verificar que el estado esté completo y persistido antes de emitir el certificado. La lección de SPARK-46182 es que una tarea terminada no implica que sus datos ya estén a salvo si queda una ventana de pérdida al cerrar \cite{spark46182}. La verificación sería por vehículo, incluso si varias líneas atienden al mismo tiempo; es una propuesta de diseño motivada por el caso de Spark y por la semántica de transacciones \cite{spark46182,postgresTransactions}.
